\section{Practical Implementation and Approximation}
\label{sec:practical}

A natural objection is that $\Sigma_2$ objects appear to require enumerating large fibers
$F(r)=\Pi^{-1}(r)$ and intervention-indexed continuation structures. This section clarifies the
intended use: $\Sigma_2$ is not a promise of universal tractability; it is a specification of
\emph{structurally necessary information}, which enables principled approximation schemes.

\subsection{Finite parameterizations of fibers}
\label{sec:finite-param}

In many applications, the physically relevant variability within $F(r)$ is low-dimensional.
For example, when the dominant uncertainty is a small set of correlation parameters (e.g.\
effective bath-mode phases, coupling strengths, or a finite set of initial correlation
coefficients), one can represent a fiber by a parameter family
\[
  F(r) \approx \{ h(\theta) : \theta \in \Theta_r \subset \mathbb{R}^k \},
\]
and define continuation structure by propagating $\theta$-families under interventions.

\subsection{Sample-based approximation of continuation sets}
\label{sec:sampling}

When parametric structure is unknown or high-dimensional, one can approximate $\Sigma_2$ by
sampling histories consistent with a reduced description:
\[
  \widehat{F}(r) = \{ h^{(1)},\dots,h^{(N)} \} \subseteq \Pi^{-1}(r),
\]
then estimating continuation sets
\[
  \widehat{A}_I(r) = \{ \Pi(h') : \exists h\in \widehat{F}(r)\ \text{s.t.}\ h \sim_I h' \}.
\]
This yields empirical set-valued predictions: interventions map uncertainty in fibers to
uncertainty in outcomes.

\subsection{Coarse-graining by continuation equivalence}
\label{sec:coarse-grain}

A particularly useful reduction is to partition $F(r)$ into equivalence classes that share
continuation structure under a restricted intervention class $\mathcal{I}_0 \subseteq
\mathcal{I}$:
\[
  h \equiv_{\mathcal{I}_0} \tilde{h}
  \quad \Longleftrightarrow \quad
  \mathrm{Cont}(\Pi(h);\mathcal{I}_0)=\mathrm{Cont}(\Pi(\tilde{h});\mathcal{I}_0).
\]
The quotient $F(r)/\!\equiv_{\mathcal{I}_0}$ is a minimal coarse representation sufficient for
counterfactual prediction under $\mathcal{I}_0$.

\subsection{Detecting when \texorpdfstring{$\Sigma_1$}{Sigma-1} is adequate}
\label{sec:detect}

Although $\Sigma_1$ descriptions can fail generically - in the structural sense clarified above - under the conditions of Proposition~\ref{prop:unfaithful-from-correlations}, there are important regimes where it is adequate in
practice (e.g.\ engineered initial factorization, restricted intervention classes, weak coupling
and short times). Operationally, one can test for unfaithfulness by designing \emph{continuation
diagnostics}: prepare multiple realizations that agree in $r$ and apply a chosen intervention
family; statistically significant multi-modal outcome structure indicates fiber sensitivity.

\subsection{What \texorpdfstring{$\Sigma_2$}{Sigma-2} changes in practice}
\label{sec:practical-summary}

In applied settings, the principal role of $\Sigma_2$ is to:
\begin{itemize}
  \item identify the structural failure mode (degenerate reductions under correlation-sensitive
        interventions),
  \item provide a target object for approximation (fiber + continuation structure),
  \item support robust design goals (controls and inference procedures that remain valid over
        fibers or fiber classes).
\end{itemize}

\subsection{Experimental detection of unfaithful cuts}
\label{sec:experimental-unfaithful}

The unfaithful cut admits a direct experimental diagnostic. Consider a controlled quantum
system in which nominally identical reduced preparations $\rho_S$ may differ in correlations
with uncontrolled environmental degrees of freedom (e.g., due to preparation pathway
variations or ambient field fluctuations).

An experimental test proceeds as follows:
\begin{enumerate}
  \item Prepare an ensemble $\{h^{(1)}_k\}$ using preparation method $P_1$.
  \item Prepare an ensemble $\{h^{(2)}_k\}$ using preparation method $P_2$.
  \item Verify that both ensembles yield statistically indistinguishable reduced states
        $\rho_S$ under standard state tomography.
  \item Apply an identical control intervention $I \in \mathcal{I}$ to both ensembles.
  \item Measure the resulting outcome distributions.
\end{enumerate}

If the post-intervention outcome distributions differ beyond statistical uncertainty, this
demonstrates fiber degeneracy: the reduction $\Pi$ is unfaithful with respect to the intervention
$I$. Conversely, agreement of outcome statistics indicates that the experiment operates within
a $\Sigma_1$-adequate regime for the chosen intervention class.

This diagnostic distinguishes genuine second-order behavior from regimes in which preparation
variations are operationally irrelevant.